\documentclass[10pt,a4paper]{amsart}
\usepackage[margin=19mm]{geometry}
\usepackage{amsmath,amssymb,mathtools}
\usepackage[scr=rsfs,cal=euler]{mathalfa}
\usepackage{xfrac}
\usepackage[unicode,hidelinks,bookmarks=false]{hyperref}
\newcommand{\ie}{i.e.\ }
\newcommand{\cf}{cf.\ }
\newcommand{\resp}{respectively\ }
\newcommand{\Subsection}{\subsection}
\providecommand{\nobreakdash}{\nobreak\hbox{-}}
\setlength{\emergencystretch}{2em}
\multlinegap0pt
\usepackage{amssymb}
\usepackage{enumerate}
\usepackage{csquotes}
\usepackage{xcolor}
\usepackage{float}
\usepackage{bm}
\newtheorem{lemma}{Lemma}
\newtheorem{theorem}{Theorem}
\newtheorem{claim}{Claim}
\title{A classification of incompleteness statements}
\author{Henry Towsner, James Walsh}
\date{Source: arXiv:2409.05973v3 (CC BY 4.0)}
\begin{document}
\maketitle

\noindent\emph{Source and licence.} A classification of incompleteness statements. Version arXiv:2409.05973v3 (CC BY 4.0), distributed by arXiv under the Creative Commons Attribution 4.0 International licence, http://creativecommons.org/licenses/by/4.0/, as recorded in the arXiv licence metadata for this exact version and shown on the abstract page https://arxiv.org/abs/2409.05973v3. Full licence notice: \texttt{licenses/arxiv-2409.05973-CC-BY-4.0.txt}.\par\smallskip

\noindent\textit{Editorial scope.} Counted unit: the article's theorem dual-positive-theorem (source0.tex lines 146--148). The article's complete original proof of it is delivered with it (source0.tex lines 149--177, one \texttt{proof} environment of 29 lines). No mathematics is written, repaired or re-derived: the statement and the proof are the article's own lines, in the article's own notation, and no retained line is substituted or re-flowed.  Retained uncounted as the source-local context the proof needs: lines 136--138.  Nothing was omitted from any retained span, so the package carries no omission record.  No retained line contains a citation, so the wrapper carries no bibliography and no bibliography entry.  No retained line contains a figure environment, an image-inclusion command or drawing code, so no image is delivered and the package declares no asset dependency.  Editorial formatting only, in addition: an AMS \texttt{amsart} preamble that restates the article's own declarations and macros used by the retained lines, namely the environments \texttt{lemma}, \texttt{theorem}, \texttt{claim} on separate counters, together with the standard packages those lines need.  The retained environments print in this document as Lemma 1, Theorem 1, Claim 1 in order of appearance, whereas the article prints the same items with the numbering of its own full document.  The complete native source of the article as distributed in the arXiv e-print for this exact version is bundled unchanged as \texttt{sources/165016/source0.tex}.
\begin{lemma}\label{reflection}
For $T$ extending $\mathsf{ACA}_0$, $\mathsf{RFN}_{\widehat\Gamma}(T)$ does not follow from any consistent extension of $T$ by $\Gamma$ formulas.
\end{lemma}

\begin{theorem}\label{dual-positive-theorem}
Let $\Gamma\in\{\Sigma^1_1,\Pi^1_1\}$. For any sound and arithmetically definable theory $U$, there is a $\widehat\Gamma$-sound and $\widehat\Gamma$-definable extension of $U$ that proves its own $\Gamma$-soundness.
\end{theorem}

\begin{proof}
Consider the following formulas:
\begin{flalign*}
    \varphi(x) &:=   x=\ulcorner\mathsf{RFN}_{\Gamma}(U)\urcorner \vee x =\ulcorner \neg \mathsf{RFN}_{\widehat\Gamma}(U + \mathsf{RFN}_{\Gamma}(U)) \urcorner\\
    \tau(x) & := U(x) \vee \Big( \mathsf{RFN}_{\widehat\Gamma}\big(U + \mathsf{RFN}_{\Gamma}(U) \big) \wedge \varphi(x) \Big)
\end{flalign*}

Let $T$ be the theory defined by $\tau$.

\begin{claim}
    $T$ is $\widehat\Gamma$-definable via $\tau$.
\end{claim}
By inspection.

\begin{claim}
    $T$ is $\widehat\Gamma$-sound.
\end{claim}
Since $U$ is sound, $U+\mathsf{RFN}_\Gamma(U)$ is sound, so $\mathsf{RFN}_{\widehat\Gamma}(U+\mathsf{RFN}_\Gamma(U))$ holds, and therefore externally, we see that $T$ is the theory:
$$  U + \mathsf{RFN}_{\Gamma}(U) + \neg \mathsf{RFN}_{\widehat\Gamma}(U + \mathsf{RFN}_{\Gamma}(U)).$$
In particular, $T$ has the form $U'+\neg \mathsf{RFN}_{\widehat\Gamma}(U')$ where $U'$ is sound.
Suppose that $U'+\neg \mathsf{RFN}_{\widehat\Gamma}(U')\vdash \sigma$ where $\sigma$ is false $\widehat\Gamma$. Then $U'+ \neg \sigma \vdash \mathsf{RFN}_{\widehat\Gamma}(U')$. So $\mathsf{RFN}_{\widehat\Gamma}(U')$ follows from a consistent extension of $U'$ by $\Gamma$ formulas, contradicting Lemma \ref{reflection}.

\begin{claim}
    $T\vdash \mathsf{RFN}_{\Gamma}(\tau)$.
\end{claim}
From our external characterization of $T$ we see that 
$$T\vdash \neg \mathsf{RFN}_{\widehat\Gamma}(U + \mathsf{RFN}_{\Gamma}(U)).$$
Hence $T$ proves that $\tau$ defines the theory $U$. Again, appealing to our external characterization of $T$, $T\vdash \mathsf{RFN}_{\Gamma}(U)$. Thus, $T\vdash \mathsf{RFN}_{\Gamma}(\tau)$.
\end{proof}

\end{document}
